%%%%%%%%%%%%%%%%%%%%%%%%%%%%%%%%%%%%%%%%%%%%%%%%%%%%%%%

\section{\scshape Motivation}

\subsection{\scshape The puzzle}

%%%%%%%%%%%%%%%%%%%%%%%%%%%%%%%%%%%%%%%%%%%%%%%%%%%%%%%

\begin{frame}{Motivation}

\onslide<1->{
	\textbf{Last two centuries: half the world transitioned from autocracy to democracy.}
}
		\vspace{.6cm}

\onslide<2->{
	\textbf{Successful democratizations are defined by changes to political institutions:}
}
		\vspace{.1cm}

	\begin{itemize}
		\item<2->[--] Free and fair elections, expanded voting rights, freedom of expression, etc.
	\end{itemize}

\vspace{.6cm}

\onslide<3->{

\textbf{Democratizations also change policy, the distribution of resources, and macroeconomic aggregates.}

}
		\vspace{.6cm}
				
\onslide<4->{\textbf{Financial markets tell us how wealthy investors perceive democratizations.}}

		\vspace{.6cm}

\onslide<5->{\textbf{Use this to test redistribution-based models of democratizations.}}

		\vspace{.6cm}


\begin{enumerate}
\item<6-> \textbf{Risk premia rise during democratizations.}
		\vspace{.3cm}
\item<7-> \textbf{Driven by fear over future redistribution.}
\end{enumerate}

\end{frame}

%%%%%%%%%%%%%%%%%%%%%%%%%%%%%%%%%%%%%%%%%%%%%%%%%%%%%%%

\begin{frame}{Risk premia rise during democratizations}
	\begin{center}
		\textbf{Log dividend yield around democratizations}
		\only<1>{
		\includegraphics[width=1\textwidth]{\figurepath/figure_1_dividend_yield_event_study.pdf} \\}
		\only<2>{
		\includegraphics[width=1\textwidth]{\figurepath/figure_1_dividend_yield_event_study.pdf} \\}
		\textbf{Gordon growth intuition: dividend yield (prior 12-month dividends divided by price) is discount rate less expected cashflow growth, $D/P = r-g$. Data from 90 countries over 200 years.}
	\end{center}	
		
\end{frame}

%%%%%%%%%%%%%%%%%%%%%%%%%%%%%%%%%%%%%%%%%%%%%%%%%%%%%%%

\begin{frame}{Main results}

	\begin{enumerate}
		\item<1-> \textbf{Risk premia rise during democratizations:}
		\begin{itemize}
			\vspace{.10cm}
			\item<1->[--] \textbf{Rule out other explanations:} macroeconomic risk and generic political risk cannot explain the results.
			\vspace{.10cm}
			\item<2->[--]  \textbf{Causal evidence:} majority Catholic autocracies have higher average excess returns after Catholic doctrinal shift favoring democracy (1959--1963).
		\end{itemize}				
		\vspace{.25cm}
	\item<1-> \textbf{Driven by fear over future redistribution:}
			\vspace{.10cm}
		\begin{itemize}
			\item<3->[--] \textbf{Redistribution after successful democratizations:} 24\% rise in tax revenue-GDP ratios, \tabEightThreeCE\ p.p. decline in Gini coefficients, \tabEightFourCE\ p.p. rise in labor share.
			\vspace{.10cm}
			\item<4->[--] \textbf{Other redistribution:} corruption and bribery indices fall and economic competition rises.
			\vspace{.10cm}
			\item<5->[--] \textbf{Asset pricing model with democratic transitions:} standard redistribution channels can explain the asset pricing results.
			\vspace{.10cm}			
			\item<6->[--] \textbf{There and back again:} model matches asset prices during autocratizations too.
			\end{itemize}
	\end{enumerate}
\end{frame}